% SGA 3，第 VIII 篇，英文重建底本说明。
% 范围：推论 1.4 至注 1.7.4，止于第 2 节之前。
% 权威源：Polo--Gille Exp8-8nov09.pdf, 组件内第 4--7 页（部分）,
% 完整合阅本第 620--623 页（部分）.
% 英文底本说明：以权威源为准；对照英文仅辅助措辞。

% 源页码：Exp8 组件内第 4 页; 完整合阅本第 620 页.
\medskip
\noindent\SGARefTarget{sga3:viii:corollary:1:4}\textbf{推论 1.4.}
\textit{可对角化群是自反的，因而局部可对角化群也是自反的。若 \(M\) 与 \(N\) 为通常的交换群，则自然同态}
\SGARefTarget{sga3:viii:corollary:1:4:display:natural-isomorphism:01}\[
  \operatorname{Hom}_{S\text{-}\mathrm{gr}}(M_S,N_S)
  \longrightarrow
  \operatorname{Hom}_{S\text{-}\mathrm{gr}}(D(N_S),D(M_S))
\]
\textit{为双射。}

\footnote{编者注：以下讨论作了扩充。}
\SGARefTarget{sga3:viii:corollary:1:4:discussion:01}由于\SGARefLink{sga3:viii:corollary:1:4:display:natural-isomorphism:01}{上述同构}与基扩张相容，它给出群值 \(S\)-函子的同构
\SGARefTarget{sga3:viii:eq:1:4:1}
\begin{equation}
  \tag{1}
  \operatorname{\underline{Hom}}_{S\text{-}\mathrm{gr}}(M_S,N_S)
  \xrightarrow{\ \sim\ }
  \operatorname{\underline{Hom}}_{S\text{-}\mathrm{gr}}
    (D(N_S),D(M_S)).
\end{equation}
对每个 \(S\)-预概形 \(T\)，
\[
  \operatorname{\underline{Hom}}_{S\text{-}\mathrm{gr}}(M_S,N_S)(T)
  =
  \operatorname{Hom}_{\mathrm{gr}}\bigl(M,\Gamma(N_T/T)\bigr),
\]
且由 I，\SGARefLink{sga3:i:subsection:1:8}{1.8}，\(\Gamma(N_T/T)\) 为局部常值映射 \(T\to N\) 所成的阿贝尔群。另一方面，令 \(\operatorname{Hom}_{\mathrm{gr}}(M,N)_S\) 为与通常的阿贝尔群 \(\operatorname{Hom}_{\mathrm{gr}}(M,N)\) 相联系的常值 \(S\)-群。显然有交换群值 \(S\)-函子的同态
\SGARefTarget{sga3:viii:eq:1:4:2}
\begin{equation}
  \tag{2}
  \operatorname{Hom}_{\mathrm{gr}}(M,N)_S
  \xrightarrow{\ \theta\ }
  \operatorname{\underline{Hom}}_{S\text{-}\mathrm{gr}}(M_S,N_S),
\end{equation}
它总是单态射，且在 \(M\) 为有限型时为同构。\footnote{编者注：分别用 \(F(M)\) 和 \(G(M)\) 表示 \SGARefLink{sga3:viii:eq:1:4:2}{(2)} 的左边和右边。若 \(M=M_1\oplus M_2\)，有典范同构 \(F(M)=F(M_1)\oplus F(M_2)\) 和 \(G(M)=G(M_1)\oplus G(M_2)\)。因此只须在 \(M=\mathbf Z/r\mathbf Z\)、\(r\geq0\) 时验证 \(\theta\) 为同构。此时 \(F(M)=({}_rN)_S\)，其中 \({}_rN\) 为 \(r\operatorname{id}_N\) 的核，且对每个 \(T\to S\)，同态
\[
  F(M)(T)=\Gamma({}_rN_T/T)
  \longrightarrow
  G(M)(T)={}_r\Gamma(N_T/T)
\]
为双射。}

\SGARefLink{sga3:viii:corollary:1:5}{下面推论}的 \SGARefLink{sga3:viii:corollary:1:5:item:a}{(a)} 由\SGARefLink{sga3:viii:corollary:1:4:discussion:01}{上述讨论}得出；\SGARefLink{sga3:viii:corollary:1:5:item:b}{(b)} 则由它结合 \SGARefLink{sga3:viii:paragraph:1-7}{1.7} 中回顾的下降结果得出。
\footnote{编者注：原文在 \SGARefLink{sga3:viii:corollary:1:5}{1.5} 后指出，更一般地，若 \(G\) 和 \(H\) 局部可对角化，且 \(H\) 为有限型，则 \(\operatorname{\underline{Hom}}_{S\text{-}\mathrm{gr}}(G,H)\) 可表。\SGARefLink{sga3:viii:corollary:1:5:item:b}{这一断言}已纳入 \SGARefLink{sga3:viii:corollary:1:5}{1.5}，其证明在 \SGARefLink{sga3:viii:paragraph:1-7}{1.7} 中明确给出。}

\medskip
\noindent\SGARefTarget{sga3:viii:corollary:1:5}\textbf{推论 1.5.}
\begin{enumerate}
  \item[(a)]\SGARefTarget{sga3:viii:corollary:1:5:item:a}
  \textit{设 \(M\) 和 \(N\) 为通常的交换群，且 \(M\) 为有限型。则有同构}
  \[
    \operatorname{Hom}_{\mathrm{gr}}(M,N)_S
    \xrightarrow{\ \sim\ }
    \operatorname{\underline{Hom}}_{S\text{-}\mathrm{gr}}
      (D(N_S),D(M_S));
  \]
  \textit{从而 \(\operatorname{\underline{Hom}}_{S\text{-}\mathrm{gr}}
  (D(N_S),D(M_S))\) 可表。}

  \item[(b)]\SGARefTarget{sga3:viii:corollary:1:5:item:b}
  \textit{更一般地，若 \(G\) 和 \(H\) 局部可对角化，且 \(H\) 为有限型，则 \(\operatorname{\underline{Hom}}_{S\text{-}\mathrm{gr}}(G,H)\) 可表。}
\end{enumerate}

\footnote{编者注：应增加一个陈述 1.5.1，讨论 X，\SGARefLink{sga3:x:corollary:5-10}{5.10} 和 5.11 中考虑的函子 \(\operatorname{\underline{Isom}}_{S\text{-}\mathrm{gr}}(G,H)\)。}
由推论~\SGARefLink{sga3:viii:corollary:1:5}{1.5} 得到：

\medskip
\noindent\SGARefTarget{sga3:viii:corollary:1:6}\textbf{推论 1.6.}
\textit{在 \SGARefLink{sga3:viii:corollary:1:5}{1.5} 的假设下，若 \(S\) 连通，则}
\[
  \operatorname{Hom}_{S\text{-}\mathrm{gr}}(D_S(N),D_S(M))
  \xrightarrow{\ \sim\ }
  \operatorname{Hom}_{\mathrm{gr}}(M,N)
\]
\textit{且}
\[
  \operatorname{Isom}_{S\text{-}\mathrm{gr}}(D_S(N),D_S(M))
  \xrightarrow{\ \sim\ }
  \operatorname{Isom}_{\mathrm{gr}}(M,N).
\]

\medskip
\noindent\SGAInlineRefTarget{sga3:viii:paragraph:1-7}\SGARefTarget{sga3:viii:subsection:1:7}\textbf{1.7. 可表性的下降。}
\footnote{编者注：增加\SGARefLink{sga3:viii:paragraph:1-7}{本段}，是为了明确指出 \(S\) 上某些层的可表性具有“关于 \(S\) 的局部性”。下文多次用到这一性质，原文中也曾隐含地使用它。}
这里回顾下文将频繁使用的若干下降结果。

\medskip
\noindent\SGARefTarget{sga3:viii:scholium:1:7:1}\textbf{附论 1.7.1.}
\textit{设 \(S\) 为预概形，\(X,Y,T\) 为 \(S\)-预概形。若 \((T_i)\) 为 \(T\) 的开覆盖，且 \(T_{ij}=T_i\cap T_j=T_i\times_TT_j\)，则由于给定 \(T\)-预概形态射 \(X_T\to Y_T\) 是关于 \(T\) 的局部问题，有集合的正合列}
\SGARefTarget{sga3:viii:eq:1:7:1}
\begin{equation}
  \tag{1}
  \operatorname{Hom}_T(X_T,Y_T)
  \longrightarrow
  \prod_i\operatorname{Hom}_{T_i}(X_{T_i},Y_{T_i})
  \rightrightarrows
  \prod_{i,j}\operatorname{Hom}_{T_{ij}}(X_{T_{ij}},Y_{T_{ij}}).
\end{equation}
\textit{因此 \(\operatorname{\underline{Hom}}_S(X,Y)\) 为局部 \(S\)-函子，即 \((\mathbf{Sch})_{/S}\) 上关于 Zariski 拓扑的层。}

更一般地，由 IV，\SGARefLink{sga3:iv:corollary:4:5:13}{4.5.13}，\(\operatorname{\underline{Hom}}_S(X,Y)\) 是 \((\mathbf{Sch})_{/S}\) 上关于每个比典范拓扑粗的拓扑的层，例如关于 fpqc 拓扑的层。若 \(G\) 和 \(H\) 为 \(S\) 上的群预概形，则子函子 \(\operatorname{\underline{Hom}}_{S\text{-}\mathrm{gr}}(G,H)\) 为 fpqc 层，从而更是局部函子。

\medskip
\noindent\SGARefTarget{sga3:viii:lemma:1:7:2}\textbf{引理 1.7.2.}
\footnote{编者注：亦见注 XI，\SGARefLink{sga3:xi:remark:3:4}{3.4}。}
\textit{设 \(F\) 为局部 \(S\)-函子。}
\begin{enumerate}
  \item[(i)]\SGARefTarget{sga3:viii:lemma:1:7:2:item:i}
  \textit{假设存在 \(S\) 的开覆盖 \((S_i)\)，使 \(F\) 在每个 \(S_i\) 上的限制 \(F_i=F\times_SS_i\) 可由 \(S_i\)-预概形 \(X_i\) 表示。则 \(F\) 可由 \(S\)-预概形 \(X\) 表示。}

  \item[(ii)]\SGARefTarget{sga3:viii:lemma:1:7:2:item:ii}
  \textit{假设 \(F\) 为 fpqc 层，并存在忠实平坦拟紧态射 \(S'\to S\)，使 \(F'=F\times_SS'\) 可由 \(S'\)-预概形 \(X'\) 表示。则 \(X'\) 带有关于 \(S'\to S\) 的下降数据（参见 IV，\SGARefLink{sga3:iv:definition:2:1}{2.1}）。若此下降数据有效——当 \(X'\) 在 \(S'\) 上仿射时即如此——则 \(F\) 可由 \(S\)-预概形 \(X\) 表示。}
\end{enumerate}

\emph{证明。}
\SGARefLink{sga3:viii:lemma:1:7:2:item:i}{(i)} 假设表明，\(X_i\times_SS_j\) 和 \(X_j\times_SS_i\) 都表示 \(F\) 在 \(S_{ij}=S_i\times_SS_j\) 上的限制。因此 Yoneda 引理给出唯一的 \(S_{ij}\)-预概形同构
\[
  c_{ji}:X_i\times_SS_j\xrightarrow{\ \sim\ }X_j\times_SS_i.
\]
在 \(S_{ijk}=S_i\times_SS_j\times_SS_k\) 上，得到同构所成的交换图
\SGARefTarget{sga3:viii:diagram:1:7:2:i}
\[
\begin{tikzcd}[column sep=small,row sep=large]
X_i\times_SS_j\times_SS_k
  \arrow[r,"c_{ji}\times\operatorname{id}_{S_k}"]
  \arrow[d,equal]
&
X_j\times_SS_i\times_SS_k
  \arrow[r,equal]
&
X_j\times_SS_k\times_SS_i
  \arrow[d,"c_{kj}\times\operatorname{id}_{S_i}"]
\\
X_i\times_SS_k\times_SS_j
  \arrow[r,"c_{ki}\times\operatorname{id}_{S_j}"]
&
X_k\times_SS_i\times_SS_j
  \arrow[r,equal]
&
X_k\times_SS_j\times_SS_i .
\end{tikzcd}
\]
事实上，它的所有对象都表示 \(F\) 在 \(S_{ijk}\) 上的限制。因此各 \(c_{ji}\) 满足通常的余循环关系 \(c_{kj}\circ c_{ji}=c_{ki}\)。

由此可将各 \(X_i\) 粘合成 \(S\)-预概形 \(X\)，且对每个 \(i\) 有 \(X\times_SS_i=X_i\)。对 \(S_i\) 上的每个 \(Y\)，有
\SGARefTarget{sga3:viii:eq:1:7:2:star}
\begin{equation}
  \tag{\(\ast\)}
  F(Y)=F_i(Y)
  =\operatorname{Hom}_{S_i}(Y,X\times_SS_i)
  =\operatorname{Hom}_S(Y,X)=h_X(Y).
\end{equation}
对任意 \(Y\to S\)，各 \(Y_i=Y\times_SS_i\) 构成 \(Y\) 的开覆盖；置 \(Y_{ij}=Y_i\times_YY_j=Y\times_SS_{ij}\)。由假设 \(F\) 是局部的，而 \(h_X\) 也是局部的，因为 Zariski 拓扑比典范拓扑粗。因此，根据 \SGARefLink{sga3:viii:eq:1:7:2:star}{\((\ast)\)}，\(F(Y)\) 和 \(h_X(Y)\) 分别与下列双箭头的核等同：
\[
  \prod_iF(Y_i)\rightrightarrows\prod_{i,j}F(Y_{ij}),
  \qquad
  \prod_ih_X(Y_i)\rightrightarrows\prod_{i,j}h_X(Y_{ij}).
\]
这证明了 \SGARefLink{sga3:viii:lemma:1:7:2:item:i}{(i)}。

\SGARefLink{sga3:viii:lemma:1:7:2:item:ii}{(ii)} 置 \(S''=S'\times_SS'\)，分别用 \(S''_1\) 和 \(S''_2\) 表示通过第一和第二投影看作 \(S'\)-预概形的 \(S''\)。函子 \(F''_1=F'\times_{S'}S''_1\) 和 \(F''_2=F'\times_{S'}S''_2\) 分别由 \(X''_1=X'\times_{S'}S''_1\) 和 \(X''_2=X'\times_{S'}S''_2\) 表示。由于
\[
  F''_1=F\times_SS''=F''_2,
\]
存在唯一的 \(S''\)-同构
\[
  c:X''_1\xrightarrow{\ \sim\ }X''_2.
\]
令 \(q_i\)、\(p_{ji}\) 分别为从 \(S'''=S'\times_SS'\times_SS'\) 到其第 \(i\) 个因子、到其第 \(i\) 和第 \(j\) 个因子的投影。令 \(X'''_i=X'\times_{S'}S'''_i\)，其中 \(S'''_i\) 表示通过 \(q_i\) 将 \(S'''\) 看作 \(S'\)-预概形；再令 \(p_{ji}^{*}(c):X'''_i\xrightarrow{\sim}X'''_j\) 由 \(c\) 经基变换得到。于是有图
\SGARefTarget{sga3:viii:diagram:1:7:2:ii}
\[
\begin{tikzcd}[column sep=large,row sep=large]
X'''_1
  \arrow[r,"p_{21}^{*}(c)"]
  \arrow[dr,"p_{31}^{*}(c)"']
&
X'''_2
  \arrow[d,"p_{32}^{*}(c)"]
\\
&X'''_3 .
\end{tikzcd}
\]
三个对象都表示 \(F\) 在 \(S'''\) 上的限制，故该图交换：
\[
  p_{32}^{*}(c)\circ p_{21}^{*}(c)=p_{31}^{*}(c).
\]
因此 \(c\) 为 \(X'\) 上关于 \(S'\to S\) 的下降数据（参见 IV，\SGARefLink{sga3:iv:definition:2:1}{2.1}）。

假设此下降数据有效，即存在 \(S\)-预概形 \(X\)，使 \(X'\simeq X\times_SS'\)。当 \(X'\) 在 \(S'\) 上仿射时，由 SGA~1，VIII，\SGARefLink{sga3:viii:proposition:2:1}{2.1}，此条件成立。\footnote{编者注：另一有效性判据见 X，\SGARefLink{sga3:x:lemma:5-4}{5.4}--\SGARefLink{sga3:viii:corollary:5:6}{5.6}。}此时对每个 \(Y\to S'\)，有
\SGARefTarget{sga3:viii:eq:1:7:2:double-star}
\begin{equation}
  \tag{\(\ast\ast\)}
  F(Y)=F'(Y)
  =\operatorname{Hom}_{S'}(Y,X\times_SS')
  =\operatorname{Hom}_S(Y,X)=h_X(Y).
\end{equation}
对任意 \(Y\to S\)，置 \(Y'=Y\times_SS'\) 和 \(Y''=Y'\times_YY'\simeq Y\times_SS''\)。如同 \(S'\to S\)，态射 \(Y'\to Y\) 也是忠实平坦且拟紧的，因而为 \(\mathcal M\)-有效满态射，其中 \(\mathcal M\) 为忠实平坦拟紧态射族。因此等价关系
\[
  Y'\times_YY'\rightrightarrows Y'
\]
为 \(\mathcal M\)-有效的，且其商为 \(Y\)。由假设 \(F\) 是 fpqc 层，而 \(h_X\) 也是，因为 fpqc 拓扑比典范拓扑粗。因此根据 \SGARefLink{sga3:viii:eq:1:7:2:double-star}{\((\ast\ast)\)}，\(F(Y)\) 和 \(h_X(Y)\) 分别与下列双箭头的核等同：
\[
  F(Y')\rightrightarrows F(Y'\times_YY'),
  \qquad
  h_X(Y')\rightrightarrows h_X(Y'\times_YY').
\]
这证明了 \SGARefLink{sga3:viii:lemma:1:7:2:item:ii}{(ii)}。\qed

\medskip
\noindent\SGARefTarget{sga3:viii:corollary:1:7:3}\textbf{推论 1.7.3.}
\textit{设 \(F\) 为 \((\mathbf{Sch})_{/S}\) 上的 fpqc 层。假设存在 \(S\) 的开覆盖 \((S_i)\)，并且对每个 \(i\) 有忠实平坦拟紧态射 \(S'_i\to S_i\)，使 \(F'_i=F\times_SS'_i\) 可由在 \(S'_i\) 上仿射的 \(S'_i\)-预概形 \(X'_i\) 表示。则 \(F\) 可由在 \(S\) 上仿射的 \(S\)-预概形 \(X\) 表示，且对每个 \(i\) 有 \(X\times_SS'_i=X'_i\)。若再有每个 \(X'_i\to S'_i\) 为闭浸入（分别为有限平展态射），则 \(X\to S\) 也是如此。}

第一个断言由 \SGARefLink{sga3:viii:lemma:1:7:2}{1.7.2} 得出。对第二个断言，只须验证每个态射 \(X\times_SS_i\to S_i\) 为闭浸入（分别为有限且平展）；这由 EGA~IV\(_2\)，2.7.1（分别再结合 IV\(_4\)，17.7.3）得出。

\medskip
\noindent\SGARefTarget{sga3:viii:remark:1:7:4}\textbf{注 1.7.4.}
如前所述，断言 \SGARefLink{sga3:viii:corollary:1:5:item:b}{1.5(b)} 由 \SGARefLink{sga3:viii:scholium:1:7:1}{1.7.1} 和 \SGARefLink{sga3:viii:lemma:1:7:2:item:i}{1.7.2(i)} 得出。
